%% das Papierformat zuerst
\documentclass[a4paper, 11pt]{article}
\usepackage[margin=3cm]{geometry}
\usepackage[utf8]{inputenc}
\usepackage[T1]{fontenc}
\usepackage[fleqn]{amsmath} %left aligned equations
\usepackage{hyperref} % clickable refs
\usepackage{graphicx}
\usepackage[toc, numberedsection]{glossaries}
\usepackage{float}
\usepackage{amssymb}
\usepackage{calc}
\usepackage{enumitem} 
\usepackage{url}
\usepackage{parskip}
% http://tex.stackexchange.com/questions/17730/newcommand-and-spacing
\usepackage{xspace}
\usepackage{xparse}
% 'frame' option for figures
\usepackage[export]{adjustbox}
% fancy page headers
\usepackage{fancyhdr}
\usepackage{xcolor}
\usepackage{tabularx}
\usepackage{listings}
\usepackage{hyperref}
\usepackage{xstring}
\usepackage[lined,boxed,commentsnumbered]{algorithm2e}
\usepackage{makecell}

% page style
\pagestyle{fancy}
\fancyhead[R]{\includegraphics[width=2cm]{../assets/kitlogo}}
\fancyhead[L]{\leftmark}
\fancypagestyle{plain}{
	\rhead{\includegraphics[width=2cm]{../assets/kitlogo}}
	\lhead{\leftmark}
}

% TODO: template übersetzten

\makeglossaries

%Hack for referencing labels
\makeatletter
\def\namedlabel#1#2{\begingroup
    #2%
    \def\@currentlabel{#2}%
    \phantomsection\label{#1}\endgroup
}
\makeatother
% End: Hack for referencing labels

% Glossar: alle Einträge, aber ohne extra Referenzen
% http://tex.stackexchange.com/questions/115635/glossaries-suppress-pages-when-using-glsaddall
\newcommand*{\glsgobblenumber}[1]{}
\makeatletter
\newcommand*{\glsaddnp}[2][]{
  \glsdoifexists{#2}{
    \def\@glsnumberformat{glsgobblenumber}
    \edef\@gls@counter{\csname glo@#2@counter\endcsname}
    \setkeys{glossadd}{#1}
    \@gls@saveentrycounter
    \@do@wrglossary{#2}
  }
}
\renewcommand{\glsaddallunused}[1][]{
  \edef\@glo@type{\@glo@types}
  \setkeys{glossadd}{#1}
  \forallglsentries[\@glo@type]{\@glo@entry}{
    \ifglsused{\@glo@entry}{}{
      \glsaddnp[#1]{\@glo@entry}}}
}
\makeatother

\renewcommand{\glsnamefont}[1]{\mdseries #1} % glossary entries shouldn’t be bold

% Glossar

% So sieht ein Glossar-Eintrag aus:
%
%\newglossaryentry{dijkstra}{
%  name={Dijkstra’s Algorithmus},
%  description={ein Algorithmus, um den optimalen Pfad in einem gerichteten Graphen zu finden}
%}
%\newglossaryentry{arc}{
%  name={Arc-Flags},
%  description={eine Technik, um Routenberechnung zu beschleunigen},
%  see={dijkstra}
%}
%
% Und so kann er im Dokument verwendet werden:
%
% lorem ipsum dolor sit \gls{arc}, consectetur
%
% End: Glossar

% usage: \counteditem{prefix}{refName} -> item `/prefixXX/` with label `prefix:refName` (where XX is counted in increments of 10)
\makeatletter
\newcommand{\oitem}[2]{
  % define the counter
  \@ifundefined{c@oitem#1}{\newcounter{oitem#1}}{} % initialized at 0
  \addtocounter{oitem#1}{10}
  \item[\namedlabel{#1:#2}{/#1\arabic{oitem#1}/}]
}
\makeatother

% new page after section
\let\oldsection\section
\renewcommand\section{\clearpage\oldsection}

\newcommand{\mamidscreenshot}[1]{\includegraphics[width=\textwidth,frame]{#1}}

\newcommand{\uiel}[3]{\item \textbf{"#1" #2:} #3}

\begin{document}

% place a symbol before clickable links
% this has to come *after* \begin{document} because hyperref installs a \AtBeginDocument hook that updates the ref command.
\newcommand{\refsymbol}[0]{\scalebox{0.5}{$\nearrow$}}
\let\oldref\ref
\renewcommand{\ref}[1]{\refsymbol\oldref{#1}}
\let\oldgls\gls
\renewcommand{\gls}[1]{\refsymbol\oldgls{#1}}
\let\oldGls\Gls
\renewcommand{\Gls}[1]{\refsymbol\oldGls{#1}}
\let\oldglspl\glspl
\renewcommand{\glspl}[1]{\refsymbol\oldglspl{#1}}
\let\oldGlspl\Glspl
\renewcommand{\Glspl}[1]{\refsymbol\oldGlspl{#1}}
\let\oldglslink\glslink
\renewcommand{\glslink}[2]{\refsymbol\oldglslink{#1}{#2}}
\let\oldhyperref\hyperref
\renewcommand{\hyperref}[2][notActuallyOptional]{\refsymbol\oldhyperref[#1]{#2}}
\let\oldautoref\autoref
\renewcommand{\autoref}[1]{\refsymbol\oldautoref{#1}}

\newcommand{\abbildung}[1]{\autoref{fig:#1}}
\newcommand{\mamid}{\textit{MAMID}\xspace}


%http://tex.stackexchange.com/questions/43002/how-to-preserve-the-same-parskip-in-minipage
\newlength{\currentparskip}
\newenvironment{minipageparskip}
{
        \setlength{\currentparskip}{\parskip}% save the value
	\begin{minipage}{\textwidth}% open the minipage
	\setlength{\parskip}{\currentparskip}% restore the value
} {
        \end{minipage}
}

\NewDocumentCommand{\refgo}{m}{%
  \IfHasLabel{#1}{%
    \hyperref[#1]{\codeinline{\refgosymboltext{#1}}}%
  }{%
   \codeinline{#1}%
   %\colorbox{red}{\texttt{#1}}%
  }%
  \let\refdisplaytext\undefined%
}

%Refgo with alternative text
\NewDocumentCommand{\refgoalt}{m m}{%
	\IfHasLabel{#1}{%
		\sloppy\hyperref[#1]{\codeinline{#2}}%
	}{%
	\codeinline{#1}%
	%\colorbox{red}{\texttt{#1}}%
}%
\let\refdisplaytext\undefined%
}

% used to produce symbol names relative to the current scope
\NewDocumentCommand{\refgosymboltext}{m}{%
  % nibble the scope hierarchy: .gocurpackage.gocurtype.gocurmethod%
  \StrLen{\gocurpackage.}[\gocurlen]%
  \StrLeft{#1}{\gocurlen}[\refgosymbolCurComponent]%
  \IfStrEq{\refgosymbolCurComponent}{\gocurpackage.}{%
    % we are in scope of the package  %
      \StrGobbleLeft{#1}{\gocurlen}[\gocurremainingpath]%
      %
      \StrLen{\gocurtype.}[\gocurlen]%
      \StrLeft{\gocurremainingpath}{\gocurlen}[\refgosymbolCurComponent]%
      \IfStrEq{\refgosymbolCurComponent}{\gocurtype.}{%
        % we are in class scope       %
        \StrGobbleLeft{\gocurremainingpath}{\gocurlen}[\gocurremainingpath]%
        %
        \StrLen{\gocurmethod.}[\gocurlen]%
        \StrLeft{\gocurremainingpath}{\gocurlen}[\refgosymbolCurComponent]%
        \IfStrEq{\refgosymbolCurComponent}{\gocurmethod.}{%
          %
          \StrGobbleLeft{\gocurremainingpath}{\gocurlen}[\gocurremainingpath]%
          \gocurremainingpath%
        }{%
          \gocurremainingpath%
        }%
      }{%
        \gocurremainingpath%
      }%
  }{%
    % we are not in scope of the package%
    #1%
  }%
}

%%%%%%%%%%%%%%%%%%%%%%%%%%%%%%%%%%%%%%%%%%%%%%%%%%%%%%%%%%%%%%%%%%%%%%%%%%%%%%%%
%                       Util-commands that are dirty hacks
%                       to enable the helper commands

% evaluate #2 if the label #1 exists, else #3.
\makeatletter
\newcommand{\IfHasLabel}[3]{%
  \@ifundefined{r@#1}% declaring a label xyz defines r@xyz
               {#3}%
               {#2}%
}
\makeatother

% Styling of inline 'code', e.g. types.
\definecolor{codeinline-gray}{gray}{0.95}
\NewDocumentCommand{\codeinline}{m}{%
  % Specify inline code / command
  % Markdown equivalent: `#1`
  \colorbox{codeinline-gray}{\texttt{#1}}%
}

%%%%%%%%%%%%%%%%%%%%%%%%%%%%%%%%%%%%%%%%%%%%%%%%%%%%%%%%%%%%%%%%%%%%%%%%%%%%%%%%

%%%%%%%%%%%%%%%%%%%%%%%%%%%%%%%%%%%%%%%%%%%%%%%%%%%%%%%%%%%%%%%%%%%%%%%%%%%%%%%%
%                               Helper Commands
% ... defining types, methods, paramters, etc

% redefine \gocurpackage to the package name that the
%   subsequent \, \property, \method, etc belong too
\newcommand{\gocurpackage}{none}
\newcommand{\gocurtype}{}
\newcommand{\gocurmethod}{}

% an abstraction for declaring types, i.e. \struc or \interface
% #1: the type of type, i.e Struct or Interface
% #2: the name of the type
% #3: prologue
% #4: properties. list of \property
% #5: methods. list of \method
%
\NewDocumentCommand{\gotype}{m m +m +g +g +d<>}{%
  \subsection{#1 \codeinline{#2}}\label{\gocurpackage.#2}%
  % Make the current type name available in the scope of the class.
  % Useful for making labels.
  \renewcommand{\gocurtype}[0]{#2}% 
  #3%
  \IfValueT{#5}{%
    \subsubsection*{Fields}%
    \begin{description}%
      #5% list of \property
    \end{description}%
   }%
   \IfValueT{#4}{%
      \subsubsection*{Methods}%
      \begin{description}%
        %\setlength{\itemsep}{1em}
        #4%list of \method
      \end{description}%
   }%
   \IfValueT{#6}{%
      \subsection*{Enum Values}%
      \begin{description}%
          #6% list of \goenumitem
      \end{description}%
   }%
   \renewcommand\gocurtype[0]{}%
}
\NewDocumentCommand{\jsview}{m +m}{\gotype{View}{#1View}{#2}}
\NewDocumentCommand{\jscontroller}{m +m}{\gotype{Controller}{#1Controller}{#2}}
\NewDocumentCommand{\struct}{m +m +g +g}{\gotype{Struct}{#1}{#2}{#4}{#3}}
\NewDocumentCommand{\interface}{m +m +g}{\gotype{Interface}{#1}{#2}{#3}}
\NewDocumentCommand{\goenum}{m +m +m}{\gotype{Enum}{#1}{#2}<#3>}
\NewDocumentCommand{\typealias}{m r[] +g}{%
  \subsection{Type Alias \codeinline{#1 $\equiv$ #2}}\label{\gocurpackage.#1}%
  #3
}

\NewDocumentCommand{\beginpackage}{m}{
	\renewcommand{\gocurpackage}{#1}\label{\gocurpackage}
}

\NewDocumentCommand{\reftype}{m}{\refgo{\gocurpackage.#1}}

% \property is used to declare fields of a struct
\NewDocumentCommand{\property}{m o +g}{
  \item[\codeinline{#1\hspace{0.1cm}:\hspace{0.1cm}#2}]\label{\gocurpackage.\gocurtype.#1}\hfill %(line break for the poor)
  \IfValueT{#3}{%
    \par
    #3
  }%
}
\NewDocumentCommand{\refproperty}{m}{\refgo{\gocurpackage.\gocurtype.#1}}

\NewDocumentCommand{\goenumitem}{m +m}{
  \item [\codeinline{#1}]\label{\gocurpackage.\gocurtype.#1} #2
}

% \method is used to declare a struct's method set.
\NewDocumentCommand{\method}{m o d() +m}{
\renewcommand\gocurmethod[0]{#1}%
\item[\codeinline{#1(\IfValueT{#3}{$\cdot$})\IfValueT{#2}{\hspace{0.1cm}:\hspace{0.1cm}#2}}]\label{\gocurpackage.\gocurtype.#1}\hfill %(line break for the poor)
  \IfValueT{#3}{
    \begin{description}
      #3 % should only contain \param
    \end{description}
  }
  \par
  #4
  \renewcommand\gocurmethod[0]{}%
}
\NewDocumentCommand{\refmethod}{m}{
	\refgo{\gocurpackage.\gocurtype.#1}
}

% \param is used to declare paramters of a \method -> check \method for details
\NewDocumentCommand{\param}{m o +m}{%
  \item[\codeinline{#1\hspace{0.1cm}:\hspace{0.1cm}#2}]\label{\gocurpackage.\gocurtype.\gocurmethod.#1} #3
}

\NewDocumentCommand{\refparam}{m}{
  \refgo{\gocurpackage.\gocurtype.\gocurmethod.#1}
}


%%%%%%%%%%%%%%%%%%%%%%%%%%%%%%%%%%%%%%%%%%%%%%%%%%%%%%%%%%%%%%%%%%%%%%%%%%%%%%%%

%%%%%%%%%%%%%%%%%%%%%%%%%%%%%%%%%%%%%%%%%%%%%%%%%%%%%%%%%%%%%%%%%%%%%%%%%%%%%%%%

%%%%%%%%%%%%%%%%%%%%%%%%%%%%%%%%%%%%%%%%%%%%%%%%%%%%%%%%%%%%%%%%%%%%%%%%%%%%%%%%

% alle Glossareintraege
\newacronym{gui}{GUI}{Graphical User Interface}
\newacronym{cli}{CLI}{Command Line Interface}
\newacronym{CRUD}{CRUD}{Create / Read / (Update|Modify) / Delete}
\newacronym{ICMP}{ICMP}{Internet Control Message Protocol}
\newacronym{API}{API}{Application Programming Interface}
\newacronym{JSON}{JSON}{JavaScript Object Notation}
\newacronym{HTTP}{HTTP}{HyperText Transfer Protocol}
\newacronym{LAN}{LAN}{Local Area Network}
\newacronym{WAN}{WAN}{Wide Area Network}

\newglossaryentry{cluster}{
	name={cluster},
	description={Aggregation of \glspl{host}},
	plural={clusters}
}
\newglossaryentry{MongoDB}{
	name={MongoDB},
	description={A NoSQL based database server},
	plural={MongoDB}
}
\newglossaryentry{replica set}{
	name={replica set},
	description={Multiple \gls{MongoDB} instances distributed over multiple \glspl{host} sharing a copy of the same data},
	plural={replica sets}
}
\newglossaryentry{sharding}{
	name={Sharding},
	description={A type of \gls{MongoDB} deployment where data sets are spread across multiple \glspl{host} or \glspl{replica set}. \\ More information: {\small \url{https://docs.mongodb.com/manual/sharding/}}}
}
\newglossaryentry{administrator}{
	name={administrator},
	description={The person responsible for managing the hardware \& software deployed on the \gls{cluster}. Uses \mamid for \gls{MongoDB} \gls{replica set} administration},
	plural={administrators}
}
\newglossaryentry{host}{
	name={host},
	description={Network node with a unique hostname},
	plural={hosts}
}
\newglossaryentry{slave}{
	name={slave},
	description={Program running on \gls{host} managing \gls{MongoDB} processes},
	plural={slaves}
}
\newglossaryentry{master}{
	name={master},
	description={Program responsible for managing \glspl{slave} and providing \acrshort{API} for \gls{cluster} management},
	plural={masters}
}
\newglossaryentry{inventory}{
	name={inventory},
	description={Persistently stored list of slaves and their state (see \ref{D:Inventory})},
	plural={inventories}
}
\newglossaryentry{active mode}{
	name={active mode},
	description={Marks the \gls{slave} as available to host \gls{MongoDB} processes that are be part of \glspl{replica set}.},
	plural={active modes}
}
\newglossaryentry{maintenance mode}{
	name={maintenance mode},
	description={Marks the \gls{slave} as unavailable \& inhibits reconfiguration through \gls{master}, but does not otherwise affect running \gls{MongoDB} instances on the \gls{host}.},
	plural={maintenance modes}
}
\newglossaryentry{disabled mode}{
	name={disabled mode},
	description={A \gls{slave} in this mode does not run any \gls{MongoDB} instances controlled by the \gls{master}, hence has no \gls{MongoDB} instances in any \gls{replica set}.},
	plural={disabled modes}
}
\newglossaryentry{risk group}{
	name={risk group},
	description={A set of \glspl{host} sharing a common risk of failure, e.g. a shared power supply. Modeled through sets of \glspl{slave} since a 1:1 relationship exists between hosts and slaves},
	plural={risk groups}
}
\newglossaryentry{persistent storage}{
	name={persistent storage},
	description={Storage capable of storing data between power outage},
	plural={persistent storages}
}
\newglossaryentry{volatile storage}{
	name={volatile storage},
	description={Storage incapable of storing data between process lifetime},
	plural={volatile storages}
}
\newglossaryentry{root data directory}{
	name={root data directory},
	description={Base directory holding all files of a \gls{slave} and its \gls{MongoDB} instances},
	plural={root data directories}
}
\newglossaryentry{arbiter}{
	name={arbiter},
	description={In charge of breaking ties on a \gls{replica set} election with even number of members. More information: {\small \url{https://docs.mongodb.com/manual/core/replica-set-elections/}}},
	plural={arbiters}
}
\newglossaryentry{degraded}{
	name={degraded},
	description={State of a \gls{replica set} with member in \gls{maintenance mode} or \gls{disabled mode} or otherwise reporting failure},
	plural={degraded}
}

\begin{titlepage}
\makeatletter
\begin{center}
~\\[4em]
{\Huge MAMID}\\[.8em]\huge{Monitor and Manager for In-memory Databases}\\[2em]
{\huge Design}\\[1em]
{\large\today}\\[2.5em]
{\LARGE
Niklas Fuhrberg\\
Anton Schirg\\
Christian Schwarz\\
Janis Streib\\
Bob Weinand\\[3em]}
{\Large supervised by}\\[2em]
{\LARGE
Dr Marek Szuba\\[1em]}
{\Large at}\\[1em]
{\LARGE
Karlsruhe Institute of Technology\\
SCC\\[2em]}
{\color{gray}
  \small Document Version: \newtheorem{theorem}{Theorem}


\section{master}

\renewcommand{\gocurpackage}{model}
\renewcommand{\gocurpackage}{master}

\struct{ClusterAllocator}{
  The \refstruct{ClusterAllocator} determines the layout of the cluster managed by \mamid.

  It attempts to fulfill the constraints defined through the model objects, in particular
  \begin{itemize}
    \item A \refgo{model.ReplicaSet}'s \refgo{VolatileNodeCount} \& \refgo{PersistentNodeCount}
    \item The \refgo{model.Slave}'s allowed number of Mongod instances
          (\refgo{MongodPortRangeBegin} to \refgo{MongodPortRangeEnd}).
    \item The \refgo{model.Slave.SlaveState}
    \item The configured \refgo{model.RiskGroup}s.
  \end{itemize}

  An iterative algorithm is employed to decide on a cluster layout described through
  \refgo{model.Mongod.DesiredState}s that
  \begin{itemize}
    \item attempts to fulfill the above constraints
    \item attempts an even distribution of Mongods on the different cluster hosts
    \item is a minimal change in comparison to the previous layout
  \end{itemize}

  \begin{theorem}{Idempotence of the ClusterAllocator}
    \label{theorem:idempotence_clusterallocator}
    Let $l$ be a layout of the cluster. Then $ClusterAllocator(ClusterAllocator(l)) = ClusterAllocator(l)$.
  \end{theorem}

}{
  \property{DB}[gorm.DB]{Initialized handle to the database.}
  \property{BusChannel}[chan interface{}]{Initialized channel to the application bus.} %TODO ref}
}{
  \method{LayoutCluster}{Lay out the cluster as described above.}
}

\subsubsection{Pseudocode}

The \refstruct{master.ClusterAllocator} is crucial to the stable operation of the \mamid-managed cluster.\\
Hence, it is worth defining the implementation of \refgo{master.ClusterAllocator.LayoutCluster} through pseudocode. %TODO parentheses

While studying the algorithms below, the reader should keep in mind that
\begin{itemize}
  \item changes in the cluster layout $\equiv$ change or creation of \refgo{model.Mongod.DesiredState} 
  \item the \refgo{model.Mongod.ObservedState} may change after an arbitrary amount or even never.\\
  \item changes to a ReplicaSet must not violate or further worsen the high-availibility constraints,
        in particular \refgo{model.ReplicaSet}'s \refgo{VolatileNodeCount} \& \refgo{PersistentNodeCount}\\
        $\implies$ Mongods in \refgo{model.MongodExecutionState.Recovering} are an important special-case.
\end{itemize}

\newcommand{\pluseq}{\mathrel{{+}{=}}}
\newcommand{\minuseq}{\mathrel{{-}{=}}}
\SetKwInOut{Input}{input}
\SetKwInOut{Output}{output}
\SetKwProg{Fn}{Function}{}{}
\SetKw{Invariant}{invariant}
\IncMargin{0.5em}

Invariant: Algorithm(Algorithm(x)) = Algorithm(x)

Otherwise oscillations could occur

Most interesting case: Set a slave to disabled. Then a new Mongod should be spawned and recover and when it is done the old disabled one can be deleted.

\begin{algorithm}

\caption{Count members of a Replica Set that are \& intended to be in stable state (running)}

\Input{ReplicaSet $r$}
\Output{Number of $p_e$ and $v_e$ member processes of $r$ that are fully operational and are planned to remain in that state.}
\BlankLine
\Fn{EffectiveMemberCount(r ReplicaSet)}{

$p_e, v_e = 0$

\For{$m \in \text{r.Mongods}$}{
	\If{$\text{m.ObservedState.ExecutionState} = Running$ \\ 	%	#this line evaluates to false if m.ObservedState = NULL
		$\land \text{m.DesiredState.ExecutionState} = Running$}{ %#this line evaluates to false if m.DesiredState = NULL 

		\uIf{$m.ParentSlave.PersistentStorage$}{
			$p_e \pluseq 1 $
		}\Else{
			$v_e \pluseq 1 $
		}
	}
}
\Return $p_e, v_e$
}
\end{algorithm}


\begin{algorithm}
\caption{Destroy members of disabled replica sets where possible without violating p/v constraints.}
\ForEach{r in ReplicaSets}{
	
	$p_e, v_e \gets $ EffectiveRunningMembers(r)
	
	\ForEach(// same for persistent and volatile){$\textbf{x}_e \in \{p_e, v_e\}$}{
		\While{$x_e > r.\textbf{x}$}{
			\textbf{destroy} any $m \in r.Mongods$ where $m.ParentSlave$ is $\textbf{x} \land \textbf{disabled}$\;
			$\textbf{x}_e \minuseq 1$
		}
		
		\Invariant minimum number of disabled slaves are running members of $r$\;
		
		\While{$\textbf{x}_e > r.\textbf{x}$}{
			\textbf{destroy} any $m \in r.Mongods$ where $m.ParentSlave$ is $\textbf{x}$\;
			$\textbf{x}_e \minuseq 1$
		}
		
		\Invariant at most $r.\textbf{x}$ members 
	}
	
	\Invariant \textit{desired state}: at most (r.p|r.v) member processes of $r$\;
	// desired state = the state that will be deployed
	
}
\end{algorithm}

\begin{algorithm}
\caption{Count recovering and active members of a Replica Set}

\Input{ReplicaSet $r$}
\Output{Number of $p_a$ and $v_a$ member processes of $r$ that \begin{itemize}
		\item are \textbf{recovering}, i.e. soon-to-be fully operational
		\item fully operational (\textbf{running})
	\end{itemize} and should actually be in that state.}
\BlankLine
\Fn{AlreadyAddedMemberCount(r ReplicaSet)}{
	$p_a, v_a = 0$
	
	\For{$m \in \text{r.Mongods}$}{
		\If{$\text{m.ParentSlave.ConfiguredState} != Disabled$ \\
			$\land \text{m.DesiredState.ExecutionState} != NotRunning$}{ 
			
			\uIf{$m.ParentSlave.PersistentStorage$}{
				$p_a \pluseq 1 $
			}\Else{
			$v_a \pluseq 1 $
		}
	}
}
\Return $p_a, v_a$
}
\end{algorithm}

\begin{algorithm}
\caption{Spawn Mongods on under-provisioned Replica Sets respecting RiskGroup \& p/v constraints.}

$p_a, v_a \gets $ AlreadyAddedMemberCount(r)

\ForEach(// same for persistent and volatile){$\textbf{x}_a \in \{p_a, v_a\}$}{

	%# TODO need to fill the queue. only ReplicaSets which 'need' (use p_* to define what need means) are in  the queue 
	$RQ \gets$ PriorityQueue(R)\;% "relative amount of missing members")\;
	$RGSQ \gets$ map[RiskGroup]PriorityQueue(Slaves of RiskGroup)\;%, "relative amount of available Mongod ports")\;
	
	\While{$r = RQ.pop()$; $r \neq nil$}{
	
		\Invariant $r$ needs a member of type $\textbf{x}$
		
		// find $x$-slave $s$ in a risk group $g$ with no other Mongod of $r$ in $g$\;
		$candidates \gets \{g \mid g \in RGSQ.keys \land \text{r has no member in g}\}$\;
		$g \gets$ ARGMAX($g \in candidates$, RGSQ[g].PeekPriority())\;
		\uIf{$g \neq nil$}{
			$s \gets RGSQ[g].pop()$\;
			\Invariant $s$ has at least one free port\;
			spawn new Mongod $m$ on $s$ and add it to $r.Mongods$\;
			compute $MongodState$ for $m$ and set the $DesiredState$ variable\;
			
			\If{$r$ needs another member of type $\textbf{x}$}{
				$RQ.insert(r)$ // recompute priority
			}
			
			\If{$s$ has free ports}{
				$RGSQ[g].insert(s)$ // recompute priority
			}
			
			
		} \Else{
			yield error\;
			continue\;
		}
		
	
	}

}

\end{algorithm}


}
\end{center}
\makeatother
\end{titlepage}
\newpage
\tableofcontents
\newpage

\section{Introduction}

This document describes the design of the Application \mamid (\emph{Monitor and Manager for In-memory Databases}).

It thorougly explains the architecture of a \mamid deployment, the distribution of responsibilites among the various components,
individual Go structures, their methods and the relationships among them.

While the textual documentation delivers detailed information on individual components, the various UML diagrams
--- and in particular the class diagram --- clarify the interoperation between components and the architecture of \mamid.

In order to facilitate lecture of this document, the reader should be familiar with the contents of the \emph{Functional Specification}
document.\\
In particular, the following sections of the \emph{Functional Specification} are crucial for proper understanding of the
application design:
\begin{itemize}
  \item the technical terms introduced listed in the \emph{Glossary}
  \item the naming of \mamid components as described in the \emph{System Model}.
\end{itemize}

\subsection{Note: UML and Go}

The authors found it difficult to apply certain concepts of the \emph{Unified Modeling Language} to a Go application design.

The following conventions were employed during the design phase:

\begin{itemize}
        \item UML \codeinline{<<enum>>} does not exist in Go. However, it can be simulated using package constants with the \codeinline{<<enum>>} type as \textbf{shared unique prefix}. 

                \textbf{Example}: \codeinline{EnumType.Item1} maps to \codeinline{EnumTypeItem1} in the Go implementation.
        %TODO more?
\end{itemize}

The reader should further recognize the remarks on the UML modeling of the \refgo{gui} in section \ref{gui:beginsection}.

\section{Deviation from the Functional Specification}

The design phase revealed flaws in the \emph{Functional Specification} that have been addressed as follows.

\subsection{Root Data Directory (MC80, F40, Section 9.3.1)}

The \emph{Root Data Directory} is no longer specified in the master's inventory using the GUI.\\
Instead, it is passed directly to the slave via a command line flag parsed by the slave executable.

Changing the value on a slave with \refgo{msp.MongodState.Running} \codeinline{Mongods} results in a situation where either
all Mongds on the slave need to be temporarily destroyed or a special-case transition needs to be implemented.
Neither alternative is desirable.

Furthermore, changing the value is usually related to manual administrative action on the specific slave host.\\
Hence, moving the parameter to the command line flags of the slave simplifies both implementation and administrative work.

\subsection{Slave Mode \emph{unknown} \& Slave Error Handling (Section 5.3.1)}

The \emph{Functional Specification} declares four slave modes: \emph{active, maintenance, disabled, unknown}.\\
The \emph{unknown} mode would be presented to the user in case of
\begin{itemize}
  \item state mismatch
  \item monitoring errors
  \item communication errors
  \item deployment errors
\end{itemize}

The error handling architecture described in this document is more transparent and allows for actual debugging of an error.

Hence, the \emph{unknown} state is replaced by more fine grained error reporting.
Since it was a state only set by the master, this change has fairly low impact on other parts of the design.

\subsection{Other Optional Criteria}

The application architecture presented in this document is designed with extensibility in mind.\\
The authors are confident that a majority of the optional criteria listed in the \emph{Functional Specification} can be implemented
transparently, i.e. without greater changes to the existing architecture.

However, the authors consider the implementation of optional criteria mostly a question of time constraints during the
implementation phase, not one of general feasibility.

\section{Overview}

\mamid is a manager for database clusters, facilitating creation, administration and monitoring of a MongoDB Replica Set deployment.

Explicit support for volatile storage on primary Replica Set members with lower-prioritized secondaries on persistent
storage is a key differentiator of \mamid.

As the manager of such distributed systems with high-availibility requirements, \mamid faces a series of non-trivial problems:
\begin{itemize}
  \item unreliable hardware
  \item unreliable communication between nodes
  \item resilience against unexpected failures of the above 
  \item finding a deployment layout that fulfills a set of availibility requirements
  \item monitoring \& state tracking of the deployed MongoDB processes
  \item abstraction of the above complexity from the administrator
  %TODO more?
\end{itemize}

In order to achieve these goals, a distributed architecture is necessary.

\mamid implements several well-established design patterns, facilitating understanding of design and implementation and improving
testability of individual components.

As specified in the \emph{Functional Specification}, the language of choice for all components but the GUI is \emph{Golang}.
While \emph{Golang} can be considered a object-oriented language, it differs from languages like \emph{Java} in certain critical approaches:
\begin{itemize}
  \item no generics
  \item no classes, only \codeinline{struct} with \emph{method sets}
  \item this is most notable in the lack of \emph{Constructors} and \emph{inheritance}
  \item however, inhertance can be simulated to a certain degree through \emph{embedding} %TODO ref go spec
\end{itemize}

Overall, this leads to a situation where many traditional design patterns do not apply as well (as e.g. to Java) and hence need
to be adopted.

Given these constraints, \mamid still employs a variety of design patterns across its different components:

% GUI Design pattersn
The architecture of the \refgo{gui} follows the \textbf{Model-View-Controller} pattern.
% is this actually MVC? @janis should elaborate on this.
% TODO explain in short the idea of MVC proposed by Angular. It should be written down somewhere -> quote it

% Master Design Patterns
The \textbf{Client/server architecture pattern} is used to decouple the \refgo{gui} functionality from the \refgo{master}:
The \refgo{gui} acts as a frontend to the functionality provided by the \refgo{master} HTTP API (\refgo{masterapi}).

Being the most complex package of the project, several decoupling patterns are employed in the \refgo{master}:

The \textbf{repository pattern} is the most one employed inside the \refgo{master}:
The external \refgo{gorm} database abstraction layer holds several tables of the structures in \refgo{model}.\\
\refgoalt{master.Monitor}{Monitor},  \refgoalt{master.ProblemManager}{ProblemManager},
\refgoalt{master.Deployer}{Deployer} and  \refgoalt{master.ClusterAllocator}{ClusterAllocator} are loosely coupled and communicate mostly by
modifying the database.\\
Furthermore, a variation of the \textbf{publish/subsribe pattern} is implemented by the \refgo{master.Bus}.

% MasterSlave Protocol
The \textbf{remote proxy} decoupling pattern is employed in the {Master Slave Protocol (\refgo{msp}) package:
it implements communication bewteen \refgo{master} and \refgo{slave}.
\refgo{msp.MSPClient} exposes RPC stubs that call the the \refgo{msp.MSPConsumer} implemented on the \refgo{slave}.
Another term used to describe this relationship between master and slaves is the \textbf{Multi-Server/Single-Client architecture}.
%TODO or is it peer2peer? check SWT slides

% Slave design patterns
\textbf{Dependency injection} is employed extensively in both \refgo{master} and \refgo{slave} to increase testability through mocks.
A good example for this are the references to \codeinline{gorm.DB} passed into \refgoalt{master.Monitor}{Monitor},  \refgoalt{master.ProblemManager}{ProblemManager},
 \refgoalt{master.Deployer}{Deployer} and  \refgoalt{master.ClusterAllocator}{ClusterAllocator}.

% Notification Manager patterns
The \textbf{strategy pattern} is used inside the \refgo{notifier} to encapsulate the concrete notification delivery
to per communication channel (\refgo{notifier.Notifier}, increasing extensibility for different communication channels in the future.
% TODO actually a strategy or more of a template method?

%TODO more design patterns

\section {Formatting \& Legend}\label{legend}

% NOTE Declaring the current package
\renewcommand{\gocurpackage}{legend}

This is a legend package explaining the formatting of the following packages.

\struct{AStruct}{
    Overall description of a \refgo{legend.AStruct}.
    It implements \refgo{legend.AnInterface}.
}{
    \property{variable}[Type]{The description of the struct's field \emph{variable}}
    \property{fieldtwo}[Type2]{
      The description of the struct's field \refgo{legend.AStruct.fieldtwo}.
    }
}{
    \method{Methodname}[ReturnType]{Method without parameters}
    \method{Methodname2}[ReturnType](%% NOTE the parentheses ()
        \param{pname}[ParamType]{Description of the parameter}
    ){
        Description of the method \refgo{legend.AStruct.Methodname2}.
        The dot in the parentheses symbolizes that this method expects parameters.
    }
}

\interface{AnInterface}{
  Overall description of the Interface. An interface has no fields, hence there is no section for fields.

  References to other types, e.g. \refgo{legend.AStruct} are possible.
  When referencing another package the package name is shown, e.g. \refgo{master.ClusterAllocator}
}{
  \method{Methodname}[ReturnType]{Method without parameters}
  \method{Methodname2}[ReturnType]( %% NOTE the parentheses ()
    \param{pname}[ParamType]{Description of the parameter}
  ){
    Description of the method \refgo{legend.AnInterface.Methodname2}.
  }
}


\section{MasterSlaveProtocol}
\renewcommand{\gocurpackage}{msp}

\begin{figure}[H]
	\includegraphics[width=\textwidth]{msp_classes}
	\caption{MasterSlaveProtocol}
\end{figure}

The MasterSlaveProtocol \refgo{msp} implements communication between \refgo{slave} and \refgo{master}.\\
There are two structures which may be sent directly over HTTP as request body, represented with JSON:
\begin{itemize}
  \item \codeinline{[]\reftype{Mongod}} An array of descriptions of all Mongod instances of that \refgo{slave}.
  \item \codeinline{\reftype{SlaveError}} Sent (only) by \refgo{slave} in case there are general problems preventing observation of any running Mongod instances.
\end{itemize}

For several datastructures in \refgo{msp}, similar ones exist in  \refgo{model}. The reasoning behind this is the desire to have a
stable protocol format for the \refgo{msp}: the JSON objects are auto-generated during runtime via reflection.
Hence, duplication of datastructures is a necessary tradeoff in order to sustain the principle of information hiding.

% server (slave) side structures
\struct{Listener}{
  Structure for handling communication through the MasterSlaveProtocol on the \refgo{slave} side.
}{
  \property{Consumer}[\reftype{Consumer}]{
    Reference to an event handler, typically a controller.
  }
}{
  \method{Run}{
    Start listening for connections using the MasterSlaveProtocol.
    \begin{itemize}
      \item Listen for incoming connections from the \refgo{master}.
      \item Decode / Encode the transport format.
      \item Validate incoming request format.
      \item Delegate valid requests to the \reftype{Consumer}.
    \end{itemize}
  }
}

\interface{Consumer}{
  A Consumer handles MasterSlaveProtocol requests received by an instance of \reftype{Listener}.
}{
  \method{RequestStatus}[{(\reftype{SlaveStatus}, \reftype{SlaveError})}]{
    %TODO reference request message object here -> proto specification
    Request to provide a status report of the slave's state.
    See \reftype{Mongod} for details on the data returned by this method.
  }
  \method{EstablishMongodState}[*\reftype{SlaveError}](
    \param{m}[\reftype{Mongod}]{State description of a Mongod instance that shall be established.}
  ){
    %TODO reference request message object here -> proto specification
    Request to establish the configuration state $m$ of a MongoDB process on the slave.
    May return an error in case the state could not be established.
  }
}


% client (master) side structures
\struct{Client}{
  Structure for handling communication through the MasterSlaveProtocol on the \refgo{master} side.
}{
  \item \textit{none}
}{
  \method{RequestStatus}[{([]\reftype{Mongod}, \reftype{error})}]{
    Send a status request to the \refgo{slave} to retrieve all the running \reftype{Mongod} instances. May return an instance of \reftype{error} in case overall retrieving failed.
  }
  \method{EstablishMongodState}[*\reftype{CommunicationError}](
    \param{m}[[]\reftype{Mongod}]{Array of \reftype{Mongod} instances to establish}
  ){
    Sends the array of \reftype{Mongod} instances to the \refgo{slave} which is tasked with creating and configuring the described processes.
  }
}


% protocol specific structures
\struct{Mongod}{
  Datastructure describing a Mongod instance controlled by the slave.
}{
  \property{Port}[\reftype{PortNumber}]{The port the Mongod instance listens on.}
  \property{ReplicaSetName}[string]{Identifier for the Replica Set the Mongod instance participates in.}
  \property{ShardingConfigServer}[bool]{If \codeinline{true}, the Mongod's configuration should be such that it can work in a MongoDB Sharding Config Server Replica Set.}
  \property{StatusError}[*\reftype{SlaveError}]{Set by \refgo{slave} when retrieving the status of the Mongod instance resulted in an error.}
  \property{LastEstablishStateError}[*\reftype{SlaveError}]{The last \reftype{SlaveError} created upon executing \refgo{msp.MSPConsumer.EstablishMongodState}.\\
       (\codeinline{nil} when sent by \refgo{master} or when state has successfully been established since the last establish request).}
  \property{State}[\reftype{MongodState}]{State the Mongod instance is / should be in.}
  \property{ReplicaSetMembers}[[]\reftype{HostPort}]{Members of the Mongod's Replica Set.\\
      When sent in response to \refgo{msp.MSPConsumer.RequestStatus}, \refparam{ReplicaSetMembers} is based on what is reported by the running Mongod instance.}
}

\goenum{MongodState}{
  Describes the state of a Mongod instance controlled by the slave.

}{
  \goenumitem{MongodStateDestroyed}{The Mongod instance is not running and all related data stored on the \refgo{slave} host is deleted. (only sent by \refgo{master})}
  \goenumitem{MongodStateNotRunning}{The Mongod instance is not running but related data may be stored on the \refgo{slave}.}
  \goenumitem{MongodStateRecovering}{The Mongod is running but not available for reads because data needs to be synced from other Replica Set members.
                                     This is a common case when adding new members to an existing Replica Set. (only sent by \refgo{slave})}
  \goenumitem{MongodStateRunning}{The Mongod instance is running and in sync with other Replica Set members.}
}

\struct{HostPort}{
  Tuple of $(hostname, port)$.
}{
  \property{Hostname}[string]{The unique hostname of a machine resolvable by the slave.}
  \property{Port}[PortNumber]{The TCP port on the host identified by \refproperty{Hostname}}
}

\typealias{PortNumber}[uint16]{
  A TCP port number as specified in \href{https://tools.ietf.org/html/rfc793\#section-3.1}{RFC 793 Section 3.1}.
}

\interface{error}{
  Generic interface for errors of the MasterSlaveProtocol. The \refgo{master.ProblemManager} will have to deal with the specific types of errors.
}

\struct{SlaveError}{
  Description of an error in the slave, usually corresponding to a partial error when executing a request from the master.
  See \refgo{msp.Mongod} for examples.
}{
  \property{Identifier}[string]{Unique identifier for the type of errror defined by the slave. Intended for communicating errors,
                                e.g. in a remote support situation. }
  \property{Description}[string]{Short human-readable description of the error, suitable for display in subject lines, etc.}
  \property{LongDescription}[string]{Long human-readable description of the error.\\
                                     Can contain the \codeinline{error.Error()} of one or more underlying errors or other useful
                                     information for further debugging.}
}

\struct{CommunicationError}{
  Description of an error generated on the \refgo{msp.MSPClient} side in case it cannot communicate properly with the \refgo{slave},
  e.g. when reading or writing to the socket fails or the connection times out.
}{
  \property{Message}[string]{Failure message from internal Go functions or message about the timeout.}
}

\section{Slave}
The slave is installed on the individual host nodes and started by the main service manager of the operating system during the boot process.

\subsection{Application / main()}{
  A Go application starts by initializing the Go 'main' package and then running the \textit{main()} function.
  This function performs early initialization of the slave's main datastructures:
  \begin{itemize}
    \item Parse command line flags using an external CLIFlagsParser.
    \item Initialize the MSPListener.
    \item Initialize the Controller with the MSPListener instance.
    \item Transfer control to MSPListener to wait for incoming connections.
  \end{itemize}
}

\struct{MSPListener}{
  Structure for handling communication through the MasterSlaveProtocol.
}{
  \property{consumer}{MSPConsumer}{
    Reference to an event handler, typically a controller.
  }
}{
  \method{Run}{
    Start listening for connections using the MasterSlaveProtocol.
    \begin{itemize}
      \item Listen for incoming connections from the Master.
      \item Decode / Encode the transport format.
      \item Validate incoming request format.
      \item Delegate valid requests to the MSPConsumer.
    \end{itemize}
  }
}

\interface{MSPConsumer}{
  An MSPConsumer handles MasterSlaveProtocol requests received by an instance of \ref{slave:MSPListener}.
}{
  \method{MSPStatusRequest}[SlaveStatus]{
    %TODO reference request message object here -> proto specification
    Request to provide a status report of the slave's state.
    See \ref{msp:SlaveStatus} for details on the data returned by this method.
  }
  \method{MSPEstablishMongodState}[error](
    \param{m}{MSP.Mongod}{State description of a mongod instance that shall be established}
  ){
    %TODO reference request message object here -> proto specification
    Request to establish the configuration state $m$ of a MongoDB process on the slave.
    May return an error in case the state could not be established.
  }
}

\struct{Controller}{
  Handles MSP requests by implementing the MSPConsumer interface. %TODO ref
  Hence, it coordinates the work required to fulfill requests from the master.
  Most importantly, it leverages ProcessManager and MongodConfigurator to spawn and configure instances of MongoDB.
}{
    \property{BusyTable}{map[PortNumber]bool}{
    Contains a boolean busy state per instance of MongoDB to prevent concurrent configuration of a single instance.
  }
  \property{busyTableLock}{sync.Mutex}{
    Mutex controlling access to busyTable. %TODO ref
  }
}

\struct{ProcessManager}{
  Starts processes, may provide a list of alive processes and eventually kill these.
}{
  \property{runningProcesses}{map[PortNumber]*exec.Cmd}{
    Holds the active process controls per instance
  }
}{
  \method{spawnProcess}[error](
    \param{m}[MSP.MongoD]{Mongod information about what exactly to spawn}
    \param{dataDir}[string]{Root directory of MongoDB data}
  ){
    Spawns a Mongod process as requested by the given Mongod inside the data root directory
  }
  \method{runningProcesses}[[]PortNumber]{
    Returns the PortNumbers of the currently running processes
  }
  \method{killProcess}[error](
    \param{p}[PortNumber]{PortNumber to identify process}
  ){
    Kills process by given PortNumber. Does not error if the process is already dead, only if it could not be killed.
  }
  \method{killProcesses}[error]{
    Kills all remaining Processes. Errors if it could not kill some processes.
  }
}

\interface{MongodConfigurator}{
  Applies or returns configuration of a mongod instance by PortNumber.
}{
  \method{MongodConfiguration}[{(MSP.MongodConfiguration, error)}](
    \param{p}{PortNumber}{PortNumber to connect to locally}
  ){
    Reads configuration from a local MongoDB instance
  }
  \method{ApplyMongodConfiguration}[error](
    \param{p}{PortNumber}{PortNumber to connect to locally}
    \param{c}{MSP.MongodConfiguration}{Configuration to apply}
  ){
    Connects to a local MongoDB instance on the given port and applies the configuration
  }
}

\struct{ConcreteMongodConfigurator}{
  Implements MongodConfigurator interface and uses an external MongoDB client to communicate with local instances
}
\newtheorem{theorem}{Theorem}


\section{master}

\renewcommand{\gocurpackage}{model}
\renewcommand{\gocurpackage}{master}

\struct{ClusterAllocator}{
  The \refstruct{ClusterAllocator} determines the layout of the cluster managed by \mamid.

  It attempts to fulfill the constraints defined through the model objects, in particular
  \begin{itemize}
    \item A \refgo{model.ReplicaSet}'s \refgo{VolatileNodeCount} \& \refgo{PersistentNodeCount}
    \item The \refgo{model.Slave}'s allowed number of Mongod instances
          (\refgo{MongodPortRangeBegin} to \refgo{MongodPortRangeEnd}).
    \item The \refgo{model.Slave.SlaveState}
    \item The configured \refgo{model.RiskGroup}s.
  \end{itemize}

  An iterative algorithm is employed to decide on a cluster layout described through
  \refgo{model.Mongod.DesiredState}s that
  \begin{itemize}
    \item attempts to fulfill the above constraints
    \item attempts an even distribution of Mongods on the different cluster hosts
    \item is a minimal change in comparison to the previous layout
  \end{itemize}

  \begin{theorem}{Idempotence of the ClusterAllocator}
    \label{theorem:idempotence_clusterallocator}
    Let $l$ be a layout of the cluster. Then $ClusterAllocator(ClusterAllocator(l)) = ClusterAllocator(l)$.
  \end{theorem}

}{
  \property{DB}[gorm.DB]{Initialized handle to the database.}
  \property{BusChannel}[chan interface{}]{Initialized channel to the application bus.} %TODO ref}
}{
  \method{LayoutCluster}{Lay out the cluster as described above.}
}

\subsubsection{Pseudocode}

The \refstruct{master.ClusterAllocator} is crucial to the stable operation of the \mamid-managed cluster.\\
Hence, it is worth defining the implementation of \refgo{master.ClusterAllocator.LayoutCluster} through pseudocode. %TODO parentheses

While studying the algorithms below, the reader should keep in mind that
\begin{itemize}
  \item changes in the cluster layout $\equiv$ change or creation of \refgo{model.Mongod.DesiredState} 
  \item the \refgo{model.Mongod.ObservedState} may change after an arbitrary amount or even never.\\
  \item changes to a ReplicaSet must not violate or further worsen the high-availibility constraints,
        in particular \refgo{model.ReplicaSet}'s \refgo{VolatileNodeCount} \& \refgo{PersistentNodeCount}\\
        $\implies$ Mongods in \refgo{model.MongodExecutionState.Recovering} are an important special-case.
\end{itemize}

\newcommand{\pluseq}{\mathrel{{+}{=}}}
\newcommand{\minuseq}{\mathrel{{-}{=}}}
\SetKwInOut{Input}{input}
\SetKwInOut{Output}{output}
\SetKwProg{Fn}{Function}{}{}
\SetKw{Invariant}{invariant}
\IncMargin{0.5em}

Invariant: Algorithm(Algorithm(x)) = Algorithm(x)

Otherwise oscillations could occur

Most interesting case: Set a slave to disabled. Then a new Mongod should be spawned and recover and when it is done the old disabled one can be deleted.

\begin{algorithm}

\caption{Count members of a Replica Set that are \& intended to be in stable state (running)}

\Input{ReplicaSet $r$}
\Output{Number of $p_e$ and $v_e$ member processes of $r$ that are fully operational and are planned to remain in that state.}
\BlankLine
\Fn{EffectiveMemberCount(r ReplicaSet)}{

$p_e, v_e = 0$

\For{$m \in \text{r.Mongods}$}{
	\If{$\text{m.ObservedState.ExecutionState} = Running$ \\ 	%	#this line evaluates to false if m.ObservedState = NULL
		$\land \text{m.DesiredState.ExecutionState} = Running$}{ %#this line evaluates to false if m.DesiredState = NULL 

		\uIf{$m.ParentSlave.PersistentStorage$}{
			$p_e \pluseq 1 $
		}\Else{
			$v_e \pluseq 1 $
		}
	}
}
\Return $p_e, v_e$
}
\end{algorithm}


\begin{algorithm}
\caption{Destroy members of disabled replica sets where possible without violating p/v constraints.}
\ForEach{r in ReplicaSets}{
	
	$p_e, v_e \gets $ EffectiveRunningMembers(r)
	
	\ForEach(// same for persistent and volatile){$\textbf{x}_e \in \{p_e, v_e\}$}{
		\While{$x_e > r.\textbf{x}$}{
			\textbf{destroy} any $m \in r.Mongods$ where $m.ParentSlave$ is $\textbf{x} \land \textbf{disabled}$\;
			$\textbf{x}_e \minuseq 1$
		}
		
		\Invariant minimum number of disabled slaves are running members of $r$\;
		
		\While{$\textbf{x}_e > r.\textbf{x}$}{
			\textbf{destroy} any $m \in r.Mongods$ where $m.ParentSlave$ is $\textbf{x}$\;
			$\textbf{x}_e \minuseq 1$
		}
		
		\Invariant at most $r.\textbf{x}$ members 
	}
	
	\Invariant \textit{desired state}: at most (r.p|r.v) member processes of $r$\;
	// desired state = the state that will be deployed
	
}
\end{algorithm}

\begin{algorithm}
\caption{Count recovering and active members of a Replica Set}

\Input{ReplicaSet $r$}
\Output{Number of $p_a$ and $v_a$ member processes of $r$ that \begin{itemize}
		\item are \textbf{recovering}, i.e. soon-to-be fully operational
		\item fully operational (\textbf{running})
	\end{itemize} and should actually be in that state.}
\BlankLine
\Fn{AlreadyAddedMemberCount(r ReplicaSet)}{
	$p_a, v_a = 0$
	
	\For{$m \in \text{r.Mongods}$}{
		\If{$\text{m.ParentSlave.ConfiguredState} != Disabled$ \\
			$\land \text{m.DesiredState.ExecutionState} != NotRunning$}{ 
			
			\uIf{$m.ParentSlave.PersistentStorage$}{
				$p_a \pluseq 1 $
			}\Else{
			$v_a \pluseq 1 $
		}
	}
}
\Return $p_a, v_a$
}
\end{algorithm}

\begin{algorithm}
\caption{Spawn Mongods on under-provisioned Replica Sets respecting RiskGroup \& p/v constraints.}

$p_a, v_a \gets $ AlreadyAddedMemberCount(r)

\ForEach(// same for persistent and volatile){$\textbf{x}_a \in \{p_a, v_a\}$}{

	%# TODO need to fill the queue. only ReplicaSets which 'need' (use p_* to define what need means) are in  the queue 
	$RQ \gets$ PriorityQueue(R)\;% "relative amount of missing members")\;
	$RGSQ \gets$ map[RiskGroup]PriorityQueue(Slaves of RiskGroup)\;%, "relative amount of available Mongod ports")\;
	
	\While{$r = RQ.pop()$; $r \neq nil$}{
	
		\Invariant $r$ needs a member of type $\textbf{x}$
		
		// find $x$-slave $s$ in a risk group $g$ with no other Mongod of $r$ in $g$\;
		$candidates \gets \{g \mid g \in RGSQ.keys \land \text{r has no member in g}\}$\;
		$g \gets$ ARGMAX($g \in candidates$, RGSQ[g].PeekPriority())\;
		\uIf{$g \neq nil$}{
			$s \gets RGSQ[g].pop()$\;
			\Invariant $s$ has at least one free port\;
			spawn new Mongod $m$ on $s$ and add it to $r.Mongods$\;
			compute $MongodState$ for $m$ and set the $DesiredState$ variable\;
			
			\If{$r$ needs another member of type $\textbf{x}$}{
				$RQ.insert(r)$ // recompute priority
			}
			
			\If{$s$ has free ports}{
				$RGSQ[g].insert(s)$ // recompute priority
			}
			
			
		} \Else{
			yield error\;
			continue\;
		}
		
	
	}

}

\end{algorithm}



%http://tex.stackexchange.com/questions/43002/how-to-preserve-the-same-parskip-in-minipage
\newlength{\currentparskip}
\newenvironment{minipageparskip}
{\setlength{\currentparskip}{\parskip}% save the value
	\begin{minipage}{\textwidth}% open the minipage
		\setlength{\parskip}{0.8em}% restore the value
	}
	{\end{minipage}}

% Args: Method, URI, 
\newcommand{\apicall}[8][]{
\begin{minipageparskip}
	\paragraph{\framebox{\uppercase{#1}} #2 - #3}
	\begin{description}[leftmargin=!,labelwidth=\widthof{\bfseries Return Codesaa}] %aa intentioal -  more distance
		\item[Parameters]
		\if\relax\detokenize{#4}\relax
		\emph{None}
		\else
		\begin{tabularx}{\linewidth}{c|c|*1{>{\centering\arraybackslash}X}@{}}
			\textbf{Parameter} & \textbf{Data Type} & \textbf{Description}\\
			\hline
			#4
		\end{tabularx}
		\fi
		
		\item[Returns]
		\begin{description}[leftmargin=!,labelwidth=\widthof{\bfseries Status Codeaa}] %aa intentioal -  more distance
			\item[Status Code] #5
			\ifthenelse{\equal{#6}{}}
			{} % optional argument #1 is empty: skip
			{\item[Headers] \begin{tabularx}{\linewidth}{c|*1{>{\centering\arraybackslash}X}@{}}
					\textbf{HTTP Header} & \textbf{Content}\\
					\hline
					#6
			\end{tabularx}}
			\ifthenelse{\equal{#7}{}}
			{\item[Body] \emph{Empty}} % optional argument #1 is empty: skip
			{\item[Body] #7}
		\end{description}
		
		\item[Error Codes]
		\begin{tabularx}{\linewidth}{c|*1{>{\centering\arraybackslash}X}@{}}
			\textbf{HTTP Status Code} & \textbf{Reason}\\
			\hline%
\ifthenelse{\equal{#1}{put}}{400 & Invalid parameters in object}{}%
\ifthenelse{\equal{#1}{get}}{404 & Object does not exist}{}%
\ifthenelse{\equal{#1}{post}}{404 & Object does not exist\\400 & Invalid parameters in object or parameter may not be changed}{}%
\ifthenelse{\equal{#1}{delete}}{404 & Object does not exist}{}%
\if\relax\detokenize{#8}\relax
\else
\\#8
\fi
		\end{tabularx}
	\end{description}
\end{minipageparskip}
\vspace{1em}
}

\section{Master API}
%TODO General api description
%Uses RESTful HTTP
%Generic error responses
\subsection{API Objects}
API Objects are exchanged in JSON format.
\subsubsection{Slave}
\begin{lstlisting}
{
	"id": <int>,
	"hostname": <string>,
	"slave_port": <int>,
	"mongod_port_range_begin": <int (inclusive)>,
	"mongod_port_range_end": <int (exclusive)>,
	"persistant_storage": <bool>,
	"root data directory": <string>,
	"state": <string>
}
\end{lstlisting}
\subsubsection{Replica Set}
\begin{lstlisting}
{
	"id": <int>,
	"name": <string>,
	"persistent_node_count": <int>,
	"volatile_node_count": <int>,
	"configure_as_sharding_config_server": <bool>
}
\end{lstlisting}
\subsubsection{Risk Group}
\begin{lstlisting}
{
	"id": <int>,
	"name": <string>
}
\end{lstlisting}
\subsection{API Methods}
\subsubsection{Slaves}
%Do not add default error codes
\apicall[get\relax]{/slaves}{Gets all slaves}
	{}
	{200}
	{}
	{array of slave objects}
	{}

\apicall[put]{/slave}{Adds a slave}
	{slave & slave object & slave to add}
	{201}
	{Location & Reference to added slave}
	{}
	{}
	
\apicall[get]{/slave/<id>}{Gets a slave}
	{id & integer & id of the slave to get}
	{200}
	{}
	{slave object}
	{}
	
\apicall[post]{/slave/<id>}{Updates a slave}
	{id & integer & id of the slave to update}
	{200}
	{}
	{}
	{}
	
\apicall[delete]{/slave/<id>}{Deletes a slave}
	{id & integer & id of the slave to delete}
	{200}
	{}
	{}
	{}
	
\subsubsection{Replica Sets}
%Do not add default error codes
\apicall[get\relax]{/replicasets}{Gets all replica sets}
	{}
	{200}
	{}
	{array of replica set objects}
	{}

\apicall[put]{/replicaset}{Adds a replica set}
	{replicaset & replica set object & replica set to add}
	{201}
	{Location & Reference to added replica set}
	{}
	{}

\apicall[get]{/replicaset/<id>}{Gets a replica set}
	{id & integer & id of the replica set to get}
	{200}
	{}
	{replica set object}
	{}

\apicall[post]{/replicaset/<id>}{Updates a replica set}
	{id & integer & id of the replica set to update}
	{200}
	{}
	{}
	{}

\apicall[delete]{/replicaset/<id>}{Deletes a replica set}
	{id & integer & id of the replica set to delete}
	{200}
	{}
	{}
	{}

\apicall[get]{/replicaset/<id>/slaves}{Gets a replica sets slaves}
	{id & integer & id of the replica set to get the slaves of}
	{200}
	{}
	{array of slave objects}
	{}

\subsubsection{Risk Groups}
%Do not add default error codes
\apicall[get\relax]{/riskgroups}{Gets all risk groups}
	{}
	{200}
	{}
	{array of risk groups objects}
	{}

\apicall[put]{/riskgroup}{Adds a risk group}
	{risk group & risk group object & risk group to add}
	{201}
	{Location & Reference to added risk group}
	{}
	{}

\apicall[get]{/riskgroup/<id>}{Gets a risk group}
	{id & integer & id of the risk group to get}
	{200}
	{}
	{risk group object}
	{}

\apicall[post]{/riskgroup/<id>}{Updates a risk group}
	{id & integer & id of the risk group to update}
	{200}
	{}
	{}
	{}

\apicall[delete]{/riskgroup/<id>}{Deletes a risk group}
	{id & integer & id of the risk group to delete}
	{200}
	{}
	{}
	{}
	
\apicall[put]{/riskgroup/<risk\_group\_id>/slave/<slave\_id>}{Assigns a slave to a risk group}
	{risk\_group\_id & integer & id of the risk group to assign the slave to\\
	 slave\_id & integer & id of the slave to assign to the risk group}
	{200}
	{}
	{}
	{}

\apicall[delete]{/riskgroup/<risk\_group\_id>/slave/<slave\_id>}{Removes a slave from a risk group}
	{risk\_group\_id & integer & id of the risk group to remove the the slave from\\
	slave\_id & integer & id of the slave to remove from to the risk group}
	{200}
	{}
	{}
	{}
\section{GUI}
\beginpackage{gui}
\begin{figure}[H]
	\includegraphics[width=\textwidth]{gui_classes}
	\caption{GUI}
\end{figure}
The \refgo{gui} is built with AngularJS. The model is defined by the returned JSON object structure of the master API (see \ref{masterapi:apiobjects}).
%TODO remark about AngularJS in the beginning of the package description

\jsview{Main}{Displays the main page.}

\jsview{ReplicaSets}{Displays all configured Replica Sets.}

\jsview{ReplicaSet}{Displays a specific Replica Set and allows CRUD of a Replica Set.}

\jsview{Slaves}{Displays all Slaves managed by \mamid.}

\jsview{Slave}{Displays a specific Slave and allows CRUD of a Slave.}

\jsview{Problems}{Displays all problems in the system.}

\jsview{RiskGroups}{Displays all risk groups and allows CRUD.}

\jscontroller{main}{}

\jscontroller{slaveIndex}{
	\begin{itemize}
		\item Uses \refgo{masterapi.slaves.getAll} to get all slaves.
	\end{itemize}
}

\jscontroller{replicaSetIndex}{
		\begin{itemize}
			\item Uses \refgo{masterapi.replicasets.getAll} to get all replica sets.
		\end{itemize}
}

\jscontroller{problemIndex}{
	\begin{itemize}
		\item Uses \refgo{masterapi.problems.getAll} to get all problems.
	\end{itemize}
}

\jscontroller{slaveById}{
		\begin{itemize}
			\item Uses \refgo{masterapi.slaves.add} to add a new slave to the system.
			\item Uses \refgo{masterapi.slaves.getById} to get a specific slave.
			\item Uses \refgo{masterapi.slaves.update} to update a specific slave.
			\item Uses \refgo{masterapi.slaves.delete} to delete a specific slave.
			\item Uses \refgo{masterapi.slaves.getProblems} to get the problems of a specific slave.
		\end{itemize}
}
\jscontroller{replicaSetById}{
		\begin{itemize}
			\item Uses \refgo{masterapi.replicasets.add} to create a new Replica Set.
			\item Uses \refgo{masterapi.replicasets.getById} to get a specific Replica Set.
			\item Uses \refgo{masterapi.replicasets.update} to update a specific Replica Set.
			\item Uses \refgo{masterapi.replicasets.delete} to delete a specific Replica Set.
			\item Uses \refgo{masterapi.replicasets.getProblems} to get the problems of a specific Replica Set.
		\end{itemize}	
}

\jscontroller{riskGroupIndex}{
	\begin{itemize}
		\item Uses \refgo{masterapi.riskgroups.getAll} to get all risk groups.
		\item Uses \refgo{masterapi.riskgroups.assignSlave} to add a slave to a specific risk group.
		\item Uses \refgo{masterapi.riskgroups.removeSlave} to remove a slave from a specific risk group.
		\item Uses \refgo{masterapi.riskgroups.update} to update a specific risk group.
		\item Uses \refgo{masterapi.riskgroups.delete} to remove a specific risk group.
	}
	\end{itemize}
}


\section{Notification Manager}
\renewcommand{\gocurpackage}{notifier}

The NotificationManager (package name: \refgo{notifier}) implements notification of a predefined list of contacts when a Problem occurs.

\section{Appendix}

In addition to this document, several manuals \& guides have been created during the implementation phase:

\begin{description}
\item[Updated Design Document] in the \codeinline{doc} repository.\\
      Represents the state of the implementation at the time of writing this document.
\item[README.md] in the \codeinline{code} repository.\\
      Gives guidance on setting up a developer machine for testing and building \mamid.
\item[INSTALL.md] in the \codeinline{code} repository.\\
      Gives guidance on deploying MAMID to a cluster, including 
      \begin{itemize}
        \item description of networking configuration
        \item help on how to use the bundled Certificate Authority helper scripts
        \item configuration for both testing and production setups
      \end{itemize}
\item[TESTING.md] in the \codeinline{code} repository.\\
      Gives guidance on how to setup the local staging environment mentioned in section \ref{dev:staging}.
\item[User Manual] for the GUI, accessible through the \mamid Web Interface (\textit{Help} button in the top bar)
\end{description}


\glsaddallunused
\makeatletter
\newglossarystyle{myAltlist}{
  \glossarystyle{altlist} % base this style on altlist
  \renewcommand*{\glossaryentryfield}[5]{
  \item[\glsentryitem{##1}\glstarget{##1}{##2}]
    \mbox{}\par\nobreak\@afterheading
    ##3\glspostdescription\space On page ##5.
  }
}
\makeatother

\end{document}
